%-----------------------------------------------------------
% Head Shell
% \label{subsec:head_shell}
%-----------------------------------------------------------

The outer shell constitutes the main structural boundary of the phantom. 
It reproduces realistic cranial curvature while maintaining manufacturability for 3D printing. 
A nominal wall thickness of \nota{t = 3.2~mm (to be confirmed)} 
was applied in the \textit{measuring region}, which corresponds to the area encompassing the 
21 electrode positions defined by the international 10–20 EEG system. 
Outside this region, the geometry was locally thickened or simplified as needed to ensure 
mechanical robustness and facilitate assembly. 
Consequently, the shell is not uniformly thick across its surface, but its 
functional behavior within the measuring area remains consistent and well controlled.

The head was divided into two parts along a \textit{superior–inferior} direction, 
resulting in an upper and a lower half rather than a mid-sagittal or planar cut. 
The separation interface was carefully positioned to remain outside the measuring region, 
thereby preventing any discontinuity or mechanical seam from affecting the probe area. 
Its geometry was shaped to ensure proper mating and alignment between both halves, 
while providing sufficient clearance to minimize potential border effects 
and to facilitate assembly and internal access.

\begin{figure}[h]
    \centering
    \includegraphics[width=0.65\textwidth]{figures/head_shell_sections.png}
    \caption{Head shell showing the two-part split and uniform wall thickness.}
    \label{fig:head_shell_sections}
\end{figure}

\nota{IMAGENES DEL CORTE DE LA CABEZA + MATING STRUCTURES PARA EL CERRADO}
